\chapter{Le procédé industriel}
\label{chap:procede}

Ce chapitre décrit le système physique, organe par organe, dans l'ordre où la matière le
traverse. Chaque section se termine par la contrainte mathématique qui en découle, dont la
formulation complète est donnée au chapitre~\ref{chap:modelisation}.

\section{Vue d'ensemble}

\begin{figure}[H]
\centering
\resizebox{\textwidth}{!}{%
\begin{tikzpicture}[
  proc/.style={rectangle, draw=ocpbleu, thick, fill=white, rounded corners=2pt,
               text width=2.9cm, align=center, minimum height=1.0cm, font=\scriptsize},
  cuve/.style={cylinder, shape border rotate=90, aspect=0.28, draw=ocpvert, thick,
               fill=ocpvert!8, text width=1.9cm, align=center, minimum height=1.0cm,
               font=\scriptsize},
  conso/.style={rectangle, draw=grisfonce, thick, fill=grisclair, rounded corners=2pt,
                text width=2.6cm, align=center, minimum height=0.9cm, font=\scriptsize},
  f/.style={-{Stealth[length=2.2mm]}, thick, draw=ocpbleu},
  boue/.style={-{Stealth[length=2.2mm]}, thick, draw=rougealerte, dashed},
]

% --- Niveau 29 ---
\node[proc] (prod29) at (0,0) {Production acide 29\\\textbf{imposée}};
\node[proc] (dec) at (0,-1.7) {Décadmiation\\$0 / 750 / 1500$};
\node[cuve] (s29s) at (3.6,0) {Stock\\29 std};
\node[cuve] (s29d) at (3.6,-1.7) {Stock\\29 dec};

% --- Niveau 54 ---
\node[proc] (conc) at (7.4,-0.85) {Concentration\\par échelons\\\textbf{volume imposé}};
\node[proc] (coc) at (11.2,0.35) {Cocristallisation\\\textbf{systématique}\\$80/20$};
\node[proc] (clar) at (11.2,-2.1) {Clarification\\décanteurs\\$90/10$};
\node[cuve] (s54) at (7.4,-3.4) {Stock\\54 NCL};

% --- Central ---
\node[cuve] (ir11) at (14.6,0.35) {IR11\\CoC + DEC\_CL};
\node[cuve] (ir12) at (14.6,-2.1) {IR12\\CL};

% --- Consommateurs ---
\node[conso] (clients) at (14.6,-4.2) {Consommateurs};

% --- Flux principaux ---
\draw[f] (prod29) -- (s29s);
\draw[f] (prod29) -- (dec);
\draw[f] (dec) -- (s29d);
\draw[f] (s29s) -- (conc);
\draw[f] (s29d) -- (conc);
\draw[f] (conc) -- (coc);
\draw[f] (conc) -- (s54);
\draw[f] (conc.east) to[out=-20,in=180] node[above,font=\tiny,pos=0.4] {dec} (clar.west);
\draw[f] (s54) to[out=0,in=210] (clar.south west);
\draw[f] (coc) -- (ir11);
\draw[f] (clar.north east) to[out=40,in=250] node[right,font=\tiny,pos=0.3] {DEC\_CL} (ir11.south);
\draw[f] (clar) -- (ir12);
\draw[f] (ir11) -- (ir12);
\draw[f] (ir12) -- (clients);

% --- Transferts interzones ---
\draw[f] (s29s.north) to[out=60,in=120,looseness=6] node[above,font=\tiny] {transferts interzones} (s29s.north);

% --- Boucles de retour ---
\draw[boue] (coc.north) to[out=100,in=60,looseness=1.1]
   node[above,font=\tiny,pos=0.5] {boue \SI{20}{\percent}} (s29s.north east);
\draw[boue] (clar.south) to[out=-90,in=-90,looseness=0.5]
   node[below,font=\tiny,pos=0.5] {boue \SI{10}{\percent}} (s29s.south);
\draw[boue] (clients.west) to[out=180,in=-90,looseness=0.7]
   node[below,font=\tiny,pos=0.75] {retour EMAPHOS \SI{40}{\percent}} (s29s.south east);

\end{tikzpicture}%
}
\caption{Schéma de flux du système. En trait plein bleu, les flux de matière descendants ;
en tirets rouges, les trois boucles de recyclage qui remontent vers le niveau 29.}
\label{fig:flux-general}
\end{figure}

La matière descend du niveau 29 vers le niveau 54, puis vers les consommateurs. Trois flux
la font \textbf{remonter} : les deux boues de traitement et le retour d'EMAPHOS. Ce sont ces
boucles qui rendent les bilans matière non triviaux.

\section{Étage 1 --- La production d'acide 29}

\subsection{Les six lignes}

Le minerai est attaqué par l'acide sulfurique ; la réaction produit de l'acide phosphorique
dilué et du gypse, séparés par filtration.

\begin{table}[H]
\centering
\caption{Production d'acide 29 du scénario de référence}
\label{tab:prod29}
\begin{tabular}{@{}lS[table-format=4.0]l@{}}
\toprule
\textbf{Ligne} & {\textbf{Production} (\tP/j)} & \textbf{Ligne 54 alimentée} \\
\midrule
13AB & 1500 & 14AB \\
13CD & 1500 & 14CD \\
13XY & 1500 & 14XY \\
13ZU & 800  & 14ZU \\
13E  & 1500 & 14EXT \\
13F  & 1500 & \textbf{aucune} \\
\midrule
\textbf{Total} & \bfseries 8300 & \\
\bottomrule
\end{tabular}
\end{table}

Deux particularités méritent d'être relevées. La ligne 13ZU produit nettement moins que les
autres : elle est partiellement à l'arrêt. Surtout, \textbf{la ligne 13F n'alimente aucune
concentration} : tout son acide part directement vers les consommateurs ou en transfert.
Elle constitue de fait la réserve souple du système, ce que le modèle retrouvera seul.

\begin{encadre}[title={Conséquence pour le modèle}]
La production est un \emph{paramètre}, noté $P_\ell$. C'est le premier point qui surprend :
dans beaucoup de problèmes académiques, la production est la variable principale. Ici, elle
est \textbf{subie}.
\end{encadre}

\subsection{La décadmiation}

Le cadmium est un métal lourd toxique naturellement présent dans le minerai. Les
réglementations européennes en plafonnent la teneur dans les engrais, d'où la nécessité de
le retirer.

Le retrait s'effectue par filtres, au niveau 29, juste après la production. \textbf{Un
filtre traite \SI{750}{\tonne} par jour} et une ligne en compte au plus deux. Il en découle
la contrainte la plus structurante de tout le modèle :

\begin{equation}
d_\ell \in \{0,\; 750,\; 1500\}
\label{eq:niveaux-decad}
\end{equation}

Ce n'est pas une variable continue. On ne peut pas décadmier \SI{312}{\tonne} : on allume
zéro, un ou deux filtres, et rien d'autre.

La décadmiation ne crée pas de matière, elle \emph{répartit} la production existante :
\begin{equation}
P_\ell = \underbrace{(P_\ell - d_\ell)}_{\text{reste standard}} + \underbrace{d_\ell}_{\text{devient décadmié}}
\end{equation}

\begin{encadre}[title={Conséquence pour le modèle}]
Le caractère discret des niveaux impose des \textbf{variables binaires}, donc une
formulation en nombres entiers. C'est ce qui fait passer le problème d'un simple programme
linéaire à un MILP.
\end{encadre}

\subsection{Les transferts interzones}

Les transferts d'acide 29 \emph{standard} entre lignes servent à rééquilibrer les stocks :
évacuer une cuve qui va déborder, alimenter une cuve qui va se vider, ou nourrir une
concentration plus gourmande que sa ligne amont.

Trois règles s'imposent.

\begin{enumerate}
  \item \textbf{Seul l'acide standard se transfère.} L'acide décadmié est traité
        immédiatement et ne circule jamais entre lignes.
  \item \textbf{Toutes les liaisons n'existent pas} (tableau~\ref{tab:matrice-interzone}).
  \item \textbf{Les quantités sont bornées} : un transfert vaut soit zéro, soit au moins
        \SI{100}{\tonne}. En deçà, l'opération n'a pas de sens industriel --- amorcer une
        pompe, mobiliser un opérateur et rincer une conduite pour quelques tonnes serait
        absurde.
\end{enumerate}

\begin{table}[H]
\centering
\caption{Matrice des transferts interzones autorisés (14 liaisons sur 30 couples possibles)}
\label{tab:matrice-interzone}
\begin{tabular}{@{}l*{6}{c}@{}}
\toprule
\textbf{De $\downarrow$ / Vers $\rightarrow$} & \textbf{E} & \textbf{F} & \textbf{AB} & \textbf{CD} & \textbf{XY} & \textbf{ZU} \\
\midrule
E  & --- & --- & $\times$ & $\times$ & --- & --- \\
F  & $\times$ & --- & $\times$ & --- & $\times$ & $\times$ \\
AB & $\times$ & --- & --- & $\times$ & --- & --- \\
CD & $\times$ & --- & $\times$ & --- & $\times$ & --- \\
XY & --- & --- & --- & $\times$ & --- & $\times$ \\
ZU & --- & --- & --- & --- & $\times$ & --- \\
\bottomrule
\end{tabular}
\end{table}

Deux observations sur cette matrice. Elle n'est \textbf{pas symétrique} : la liaison
CD~$\to$~E existe alors que E~$\to$~F n'existe pas, ce qui traduit des conduites à sens
unique ou des dénivelés. Et la ligne F est le meilleur émetteur, avec quatre destinations,
tout en ne pouvant recevoir de personne --- cohérent avec son statut de ligne
structurellement excédentaire.

\begin{encadre}[title={Conséquence pour le modèle}]
Un transfert est une variable \textbf{semi-continue} : soit nulle, soit comprise entre
\SI{100}{\tonne} et une borne supérieure. C'est la deuxième source de variables binaires.
\end{encadre}

\section{Étage 2 --- La concentration}
\label{sec:concentration}

\subsection{Le rendement de 100\,\%}

L'acide est chauffé sous vide ; l'eau s'évapore et le titre monte de \SI{26}{\percent} à
\SI{50}{\percent}. La documentation affirme un \emph{rendement massique de
\SI{100}{\percent}}, ce qui paraît d'abord contredire la conservation de la masse.

\begin{encadre}[title={La résolution de l'apparent paradoxe}]
Prenons \SI{1000}{\tonne} de solution à \SI{26}{\percent}, soit \SI{260}{\tonne} de
\PdeuxOcinq{}. Après concentration à \SI{50}{\percent}, ces \SI{260}{\tonne} sont diluées
dans \SI{520}{\tonne} de solution.

\medskip
\begin{tabular}{@{}lrcl@{}}
Masse d'acide : & \SI{1000}{\tonne} & $\to$ & \SI{520}{\tonne} \quad(\textbf{non conservée}) \\
Masse de \PdeuxOcinq{} : & \SI{260}{\tonne} & $\to$ & \SI{260}{\tonne} \quad(\textbf{conservée}) \\
\end{tabular}
\medskip

Le rendement vaut $1$ \emph{si et seulement si} l'unité de compte est la tonne de
\PdeuxOcinq{}. C'est la deuxième preuve de la section~\ref{sec:unite}.
\end{encadre}

\subsection{Les échelons}

Un échelon est une unité élémentaire de concentration. Sa production journalière vaut
\begin{equation}
\kappa_e = C_e \times \frac{h_e}{24}
\label{eq:prod-echelon}
\end{equation}
où $C_e$ est sa capacité à plein régime et $h_e$ son nombre d'heures de marche, qui est une
donnée d'entrée.

\begin{table}[H]
\centering
\caption{Répartition des échelons par ligne de concentration}
\label{tab:echelons}
\begin{tabular}{@{}llS[table-format=4.0]l@{}}
\toprule
\textbf{Ligne} & \textbf{Échelons (capacité \tP/j)} & {\textbf{Total}} & \textbf{Cocristallisation} \\
\midrule
14EXT & E(420) F(420) G(420) H(420) & 1680 & tous \\
14AB & A(300) B(300) I(590) J(590) K(300) L(300) & 2380 & sous-ensemble \\
14CD & C(250) D(250) M(290) N(250) & 1040 & aucun \\
14XY & X(300) Y(250) P(250) Q(250) & 1050 & aucun \\
14ZU & Z(250) U(300) R(250) S(300) V(590) W(590) & 2280 & aucun \\
\bottomrule
\end{tabular}
\end{table}

\subsection{La production d'acide 54 est imposée}

Les heures de marche étant une donnée, la production de chaque ligne est entièrement
déterminée avant toute optimisation :
\begin{equation}
\Pi_m = \sum_{e \in \Ech_m} \kappa_e
\end{equation}

Il en découle que la quantité d'acide 29 à envoyer en concentration n'est pas libre : elle
doit \emph{exactement} égaler ce que la ligne va produire. On ne décide pas \emph{combien}
concentrer, mais seulement quelle \textbf{proportion d'acide décadmié} mettre dans
l'alimentation.

\subsection{Le chemin imposé de l'acide décadmié}

Lorsqu'on concentre de l'acide 29 décadmié, on obtient de l'acide 54 NCL décadmié, dont le
parcours est entièrement contraint :

\begin{center}
\small
acide 29 dec $\longrightarrow$ concentration $\longrightarrow$ 54 NCL DEC
$\longrightarrow$ \textbf{clarification obligatoire} $\longrightarrow$ DEC\_CL
$\longrightarrow$ IR11
\end{center}

Il n'existe \emph{aucun} stock local pour le 54 NCL décadmié : il est clarifié
immédiatement.

\begin{encadre}[title={Une simplification élégante}]
L'entrée de clarification décadmiée est donc \emph{identiquement égale} à l'acide décadmié
concentré. Ce n'est pas une décision indépendante, mais une simple substitution. Elle
supprime cinq variables et cinq contraintes, et rend automatique le couplage entre
décadmiation et capacité de décantation.
\end{encadre}

\section{Étage 3a --- La cocristallisation}

Un refroidissement contrôlé provoque la cristallisation de l'acide ; les impuretés restent
en solution et les cristaux, très purs, sont récupérés puis refondus. Le procédé rend
\SI{80}{\percent} d'acide CoC dirigé vers IR11, et \SI{20}{\percent} de boue recyclée vers le
stock d'acide 29 standard. Il n'existe que sur les lignes 14EXT et 14AB.

\begin{alerte}[title={Le point le plus contre-intuitif du sujet}]
La cocristallisation \textbf{ne s'adapte pas à la demande}. Si un échelon lui est affecté,
\emph{tout} l'acide qu'il produit part en cocristallisation --- sans modulation, même si
personne ne demande de CoC.

Industriellement, cela s'explique : ces unités ne se démarrent et ne s'arrêtent pas
facilement, on les fait tourner en continu.
\end{alerte}

Dans le scénario de référence, cette règle produit \SI{2431}{\tonne} de CoC par jour, pour un
besoin en acide de qualité décadmiée de \SI{987}{\tonne} seulement --- un facteur $2{,}5$.

\begin{encadre}[title={Conséquence pour le modèle}]
La cocristallisation ne donne lieu à \textbf{aucune variable de décision} : elle est
entièrement calculée en prétraitement. En faire une décision produirait un modèle plus gros,
plus lent, et surtout \emph{faux}, puisqu'il proposerait des plans que l'usine ne peut pas
exécuter.
\end{encadre}

\section{Étage 3b --- La clarification}

La décantation par gravité sépare les solides en suspension. Le procédé rend
\SI{90}{\percent} d'acide clarifié et \SI{10}{\percent} de boue recyclée au niveau 29.

Deux usages se partagent la même ressource.

\begin{table}[H]
\centering
\caption{Les deux usages concurrents des décanteurs}
\label{tab:usages-decanteurs}
\begin{tabular}{@{}llll@{}}
\toprule
\textbf{Usage} & \textbf{Entrée} & \textbf{Destination} & \textbf{Nature} \\
\midrule
Clarification ordinaire & 54 NCL & IR12 & \textbf{décision} de l'optimiseur \\
Clarification décadmiée & 54 NCL DEC & IR11 & \textbf{obligatoire} \\
\bottomrule
\end{tabular}
\end{table}

Chaque ligne dispose de deux décanteurs de \SI{500}{\tonne} par jour, soit une capacité
\emph{entrante} de \SI{1000}{\tonne} par jour, tous usages confondus.

\begin{encadre}[title={Le couplage central du problème}]
Cette unique contrainte crée un arbitrage en cascade :

\begin{center}
\small
décadmier plus $\rightarrow$ plus de DEC à clarifier (obligatoire) $\rightarrow$
moins de capacité pour le NCL $\rightarrow$ moins de CL vers IR12 $\rightarrow$
risque sur la demande de U53
\end{center}

Et symétriquement, décadmier moins met en péril la qualité du mélange d'IR11. C'est cet
arbitrage que l'optimisation résout, et qu'un planificateur ne peut pas trancher de tête.
\end{encadre}

\section{Les stockages centraux}

Le bac IR11 reçoit le CoC et le DEC\_CL, \emph{mélangés} ; le bac IR12 reçoit l'acide
clarifié ordinaire. Une règle de distribution s'applique à tous les acides traités :

\begin{encadre}
\textbf{Aucun acide traité n'est expédié directement depuis sa ligne de production.} Les
qualités CL, CoC et DEC\_CL transitent obligatoirement par IR11 ou IR12.

\emph{Conséquence mathématique heureuse} : puisque les bacs centraux desservent tous les
consommateurs, ces trois qualités échappent aux contraintes d'interconnexion. Seuls les
acides livrés depuis un stock local y sont soumis.
\end{encadre}

\section{Synthèse : ce qui est décidé, ce qui est subi}

Le tableau~\ref{tab:decide-subi} résume l'ensemble du chapitre. C'est la table à retenir.

\begin{table}[H]
\centering
\caption{Frontière entre paramètres et variables}
\label{tab:decide-subi}
\begin{tabular}{@{}lp{3.1cm}p{6.3cm}@{}}
\toprule
\textbf{Élément} & \textbf{Nature} & \textbf{Justification} \\
\midrule
Production d'acide 29 & Paramètre & imposée par l'amont \\
Heures de marche & Paramètre & planning et maintenance \\
Production d'acide 54 & Paramètre & déduite des heures de marche \\
Échelons cocristallisants & Paramètre & configuration du scénario \\
Production de CoC & Paramètre & systématique, entièrement déduite \\
Rendements & Paramètre & physique du procédé \\
Matrices d'interconnexion & Paramètre & tuyauterie existante \\
Capacités & Paramètre & équipements installés \\
Demandes & Paramètre & imposées par le commerce \\
\midrule
Niveau de décadmiation & \textbf{Variable} & discrète : $\{0, 750, 1500\}$ \\
Répartition std/dec & \textbf{Variable} & continue, somme imposée \\
Quantité clarifiée & \textbf{Variable} & continue, sous capacité \\
Transferts interzones & \textbf{Variable} & semi-continue \\
Affectation ligne $\to$ client & \textbf{Variable} & continue, sous interconnexion \\
Prélèvements sur IR11/IR12 & \textbf{Variable} & continue \\
\bottomrule
\end{tabular}
\end{table}

\begin{encadre}[title={La leçon de modélisation}]
Le système paraît immense, mais l'espace de décision est petit. L'essentiel du travail
consiste à \textbf{calculer correctement les paramètres} --- production d'acide 54,
cocristallisation systématique, besoins issus des recettes --- \emph{avant} de construire le
modèle d'optimisation.

C'est exactement l'architecture retenue au chapitre~\ref{chap:implementation} :
prétraitement déterministe, puis optimisation, puis validation.
\end{encadre}
